% SJTU internship report. Single-column, first person, standalone.
%
%   pdflatex main && pdflatex main
%
% This document and paper.tex describe the same work for different audiences:
% this one is the report of what I did over the internship, including the
% benchmark construction and the mistakes I corrected; paper.tex is the
% conference/arXiv submission. Every shared number must agree between them --
% check_paper_numbers.py enforces that for paper.tex, and the tables here are
% transcribed from the same results files.

\documentclass[11pt,a4paper]{article}

\usepackage[margin=2.6cm]{geometry}
\usepackage{amsmath}
\usepackage{amssymb}
\usepackage{graphicx}
\usepackage[hidelinks]{hyperref}

% Optional packages: the document degrades gracefully without them, so it
% compiles on a minimal TeX Live as well as on a full installation.
\IfFileExists{booktabs.sty}{\usepackage{booktabs}}{%
  \newcommand{\toprule}{\hline}
  \newcommand{\midrule}{\hline}
  \newcommand{\bottomrule}{\hline}
  \newcommand{\cmidrule}[2][]{}%
}
\IfFileExists{microtype.sty}{\usepackage{microtype}}{}

% Figures are produced by make_figures.py (matplotlib) rather than TikZ, so the
% numbers in the plots come from one source and cannot drift from the text.

\newcommand{\xhat}{\hat{\mathbf{x}}}
\newcommand{\yhat}{\hat{\mathbf{y}}}
\newcommand{\zhat}{\hat{\mathbf{z}}}
\newcommand{\svec}{\mathbf{s}}
\newcommand{\gvec}{\mathbf{g}}
\newcommand{\SE}{\mathrm{SE}(3)}
\newcommand{\SO}{\mathrm{SO}}

\title{\large Internship Report}
\author{Andrea Emir Sevincel\\
        {\normalsize Shanghai Jiao Tong University}\\
        {\normalsize Supervisors: Xinzhe Zhou, Xiaoming Duan}}
\date{August 2026}

% Absorbs the last few points on paragraphs TeX cannot break well -- here,
% lines carrying PointMass3D and RRT-Connect, which resist hyphenation. Only
% applies where a line would otherwise overflow, so it does not loosen the rest.
\setlength{\emergencystretch}{1em}

% Type 1 bullets: the LaTeX default is rasterised on this TeX Live install
% and shows up as a Type 3 font, which IEEE PDF eXpress rejects.
\renewcommand{\labelitemi}{$\bullet$}

\begin{document}
\maketitle

\begin{abstract}
\noindent
This report describes what I did over the internship at Shanghai Jiao Tong University. I mainly did the following:

I built the benchmark, PointMass3D, with an expert pipeline over RRT-Connect,
CHOMP and TrajOpt; then a conditional flow-matching planner over whole
trajectories; then the symmetry study.

Varying only the representation
raises the collision-free rate from \%14.60 to
\%51.10 at convergence over three seeds, a gap of
+36.5 points, and the gap \emph{widens} with training data rather than
shrinking: across a $12.5\times$ range of environments at a fixed budget it goes
from $+6.4$ to $+30.1$. All three mechanisms built for the sixth degree of freedom are worth
almost nothing at convergence: roll augmentation $+0.8 \pm 0.3$, frame averaging
$+0.68$, and an exactly $\SO(2)$-equivariant backbone $+0.15$, seed-matched. The
architecture is not pointless, much faster (2x to 3x) than the unconstrained model, so the symmetry buys
time, not accuracy. I also bound my own contribution from above: supplying each
waypoint with the signed distance to the scene and its gradient is worth $+40.0$
points where the reduction is worth $+36.5$, so the representation change studied
here is not the largest effect available on this benchmark, and yet, with that
information given, the canonicalization is still worth $+24.8$. I explain this  by measuring  the \emph{non-equivariance residual}: only about $1.7\%$ of the learned velocity
field lies outside the equivariant subspace, and frame averaging is exactly the
orthogonal projection onto it, so that $1.7\%$ is the whole budget any post-hoc
symmetrisation of this model can spend. Two further findings support the reading:
representation and optimisation are separate bottlenecks, and the whole pattern
replicates in a second state space where a state is a full rigid-body pose.
\end{abstract}


\section{What I did, in the order in which I did it}

The internship ran in three stages, and this report follows them.

\begin{description}
  \item[benchmark.] I built PointMass3D: a cluttered 3D
  workspace with an analytic signed distance field, and an expert pipeline
  combining RRT-Connect, random shortcutting, uniform resampling and CHOMP
  refinement. Implementing three classical planners rather than one also gave me
  the \emph{ceiling} the learned planner should be read against.
  \item[the baseline planner.] I also built a conditional flow-matching model
  over whole trajectories, and the calibration that turned out to matter most:
  measuring what a straight line and an untrained stochastic prior achieve on
  the identical problems.
  \item[the symmetry study.] The $(\svec,\gvec)$ reduction, the two
  obstructions that rule out any continuous gauge for the sixth degree of
  freedom, the three mechanisms built for it, and after doing everything and finding not much, explaining why they are
  worth little at the budget of the main grid, and why that qualifier matters.
\end{description}

\subsection{Contributions}

\textbf{(i) Five of six degrees of freedom are free.} The query supplies a
frame. Placing the origin at the midpoint of $\svec$ and $\gvec$ and aligning the
first axis with the start--goal direction removes three translational and two
rotational degrees of freedom exactly, at no computational cost and with no
architectural constraint. The resulting representation is invariant under any
rigid motion of the problem \emph{up to a single roll about the start--goal
axis}.

\textbf{(ii) No continuous rule fixes the sixth, and the natural scene-derived
alternative carries no signal.}
A continuous rule fixing the roll would be a nowhere-vanishing tangent field on
$S^2$, which the hairy ball theorem forbids. Deriving the roll from the scene
instead fails because the obstacle law is symmetric under central inversion, so
the angular moment such a gauge would use vanishes in expectation.

\textbf{(iii) The symmetry machinery buys little, and I measured why.} Rather
than infer inertness from flat quality curves, I measured the quantity frame averaging
exists to remove, and proved exactly what it bounds.

\textbf{(iv) Calibration, and two bottlenecks.} A collision-free rate on a new
benchmark means nothing without a scale. Measuring the floor is what revealed
that the world-frame arm sits exactly on the straight line, and measuring a
longer training run is what revealed that compute and representation relieve
different constraints.

I regard (iii) as the transferable part. The residual is a few lines of code,
costs one forward pass, needs no retraining, and can be computed on any model
with a group action on its inputs. It converts ``we tried equivariance and it did
not help'' into a quantitative statement with an exact bound attached --- and it
comes with an equally exact statement of its limits, which
Sec.~\ref{sec:twosigns} makes concrete.

\section{Stage 1: The PointMass3D benchmark}
\label{sec:benchmark}

The workspace is $[-1,1]^3$ holding forty obstacles, twenty spheres and twenty
oriented boxes, with an analytic signed distance field so that collision
checking is exact rather than sampled. The robot is a point of radius $0.03$.
A trajectory is $N=64$ waypoints in $\mathbb{R}^3$.

Expert paths come from RRT-Connect, then random shortcutting, uniform
resampling, and CHOMP refinement. I implemented TrajOpt as well, which was not
needed for the dataset but gave me a second and third point on the ceiling the
learned planner should be read against. That turned out to matter more than the
dataset itself: a collision-free rate means nothing on a new benchmark until you
know what a trivial method scores and what a good classical method scores on the
identical problems.

The released set is $300$ environments $\times\ 36{,}000$ trajectories, from
$600$ start--goal pairs each with $30$ paths and their time reversals.
Environments $250$--$299$ are held out and never trained on. One caveat I want
on record: thirty near-duplicate paths per pair means the nominal ten million
trajectories carry on the order of $360{,}000$ distinct queries, so the dataset
is far less diverse than its size suggests. This does not affect any comparison
here, since both arms see identical data, but it does mean that if I generated
more data I would cut paths per pair and buy environments instead.

\section{Stage 2: The flow-matching planner}
\label{sec:planner}

The planner is a conditional flow-matching model over whole trajectories.
Training draws $x_0 \sim \mathcal{N}(0,I)$ and $t \sim U[0,1]$, forms
$x_t = (1-t)x_0 + t\,x_1$ with $x_1$ an expert trajectory, and regresses the
velocity $x_1 - x_0$. Sampling integrates the learned field from noise with eight
Euler steps. The backbone is a dilated temporal convolutional trunk with FiLM
conditioning and a PointNet-style obstacle encoder, $2.16$M parameters.

I considered a Brownian-bridge prior, which would start sampling from paths that
already join the endpoints, and did not adopt it. The reason I originally gave
was wrong --- I thought the bridge was not isotropic in the plane normal to the
start--goal axis, and it is --- but the decision survives for a better reason:
the bridge is query-dependent and its endpoint covariance is singular, which
complicates the comparison this report exists to make.

\section{Stage 3: Method}
\label{sec:method}

\subsection{The query already supplies a frame}

A planning problem is unchanged if you pick it up and move it: for any rigid
$T \in \SE$, the solution to the transformed problem is the transformed
solution. A network trained on world coordinates cannot see this. It must learn
the same manoeuvre again at every position and orientation.

The start and goal define a frame without any learning. Place the origin at the
midpoint of $\svec$ and $\gvec$, align the first axis with the start--goal
direction, and express every obstacle and every waypoint in that frame. Three
translational and two rotational degrees of freedom vanish exactly, at
initialisation, at no computational cost and with no architectural constraint.
As a side effect the six numbers used to condition the network on the query
collapse to a single invariant scalar, the separation $d$.

\paragraph{What survives, precisely.} This is the point I got wrong for several
weeks and it matters. The reduced representation is \emph{not} invariant under
$\SE$. The gauge needs a second axis, and mine is built from a fixed world axis,
so rotating the whole problem selects a different one and rolls the frame. What
is true is that the entire residual is that one roll: every reduced quantity ---
waypoints, sphere centres, box edge vectors --- maps through one shared rotation
about the first axis, verified to $10^{-15}$. The reduction collapses $\SE$ to
$\SO(2)$, not to nothing. Invariant outright are $d$, the first coordinate of
every reduced point, and its radius in the plane normal to the axis.

The correction makes the report more coherent, not less. If the representation
were fully invariant there would be no roll left to handle, and the whole of the
next section would be unmotivated.

\subsection{No continuous rule fixes the sixth degree of freedom}

Two independent obstructions, and I want to be careful about what each one
forbids.

\paragraph{Topological.} A rule assigning a perpendicular direction
$\yhat(\xhat)$ continuously over all axis directions is exactly a nowhere-%
vanishing tangent field on $S^2$, which the hairy ball theorem forbids. The
start--goal directions in this dataset cover the whole sphere, so there is no
contractible-domain escape. Every deterministic gauge therefore jumps somewhere;
for the smallest-component cross-product construction I use, the jump is
$60$--$120$ degrees.

\paragraph{Distributional.} One might derive the roll from the scene instead, by
pointing the second axis at the obstacle mass. That is the first angular moment
of the obstacle distribution, and it vanishes in expectation: the obstacle law is
symmetric under $\mathbf{p} \mapsto -\mathbf{p}$, which preserves every radius
and sends $\varphi \mapsto \varphi + \pi$, so every \emph{odd} moment is zero for
every axis direction. Even moments survive this argument but not the topological
one, since a second-moment gauge defines a director field and a continuous
director field on $S^2$ is obstructed for the same Euler-characteristic reason.

What these rule out is a \emph{continuous} gauge. A discontinuous rule exists and
I use one. So the sixth degree of freedom has to be handled by symmetry rather
than by canonicalisation.

\subsection{Three ways to handle the leftover roll}

A symmetry can be imposed in exactly three places, and I built all three.

\textbf{In the data.} Roll augmentation: apply a uniform random roll to each
training example, so the network sees every orientation of the residual.

\textbf{In the operator.} Frame averaging: at inference, evaluate the velocity
field at $K$ equally spaced rolls, map each result back, and average. This makes
the sampled field exactly $C_K$-equivariant regardless of what the network
learned.

\textbf{In the weights.} An $\SO(2)$-equivariant architecture, in which
equivariance holds by construction rather than by training. I describe it in
Sec.~\ref{sec:arch}.

\subsection{The non-equivariance residual, and exactly what it bounds}

Rather than infer that these mechanisms are inert from flat quality curves, I
measured the quantity frame averaging exists to remove. Let $f$ be the learned
velocity field and $\mathcal{A}f$ its average over $K$ rolls. The
\emph{non-equivariance residual} $r$ is the RMS disagreement among the rotated
copies, normalised by the field magnitude. It costs one forward pass and needs no
retraining.

$\mathcal{A}$ is linear, idempotent and self-adjoint under the natural inner
product, so it is the \emph{orthogonal projection} onto the $C_K$-equivariant
subspace. Since the target field is itself equivariant, Pythagoras gives a bound:
the error frame averaging can remove is at most $r^2$. Measured across four data
scales, $r \approx 1.7\%$ and is flat. So the entire budget available to any
post-hoc symmetrisation of this model is under $3 \times 10^{-4}$, while the
model is failing on most of its samples. \emph{The model's error is
overwhelmingly equivariant error.}

I am equally explicit about what this does not license, and
Sec.~\ref{sec:twosigns} makes it concrete: $r$ bounds what symmetrising a
\emph{given} field can remove, and does not forecast what a symmetry
\emph{mechanism} is worth.

\subsection{Implementation, and the errors worth recording}

\paragraph{Oriented boxes.} The reduction is a general rotation, so axis-aligned
boxes do not stay axis-aligned. I changed the box representation to a centre plus
three half-edge vectors, and gave both arms that representation so it is not a
confound.

\paragraph{Points and free vectors transform differently.} The reduction is
affine. Waypoints, start, goal and obstacle centres take the full affine map; the
three half-edge vectors take the rotation \emph{only}. Applying the affine map to
an edge vector leaves box positions correct and silently destroys box extent, and
the loss curve looks entirely healthy while it happens. The same class of error
later appeared in the rigid-body domain's rotation columns. It is the single
most dangerous bug I hit, because nothing fails loudly.

\paragraph{Verification.} Sixty-one tests across six files, all property-based.
Every equivariance test runs on an \emph{untrained} network, which is the case
that shows the property comes from the structure of the weights rather than from
training. Each also has a negative control: a network that ignored its vector
inputs entirely would be trivially equivariant, so the tests check that the model
is \emph{not} equivariant to rotations off the relevant axis.

\section{Protocol and calibration}
\label{sec:protocol}

Every arm is evaluated on the same $500$ held-out problems: $50$ environments
$\times$ $10$ distinct start--goal pairs $\times$ $20$ samples, eight Euler
steps. I report two metrics and never conflate them. A sample is \emph{free} if
the whole path is collision-free; \emph{best-of-20} is whether any of the twenty
samples for a query is free. They differ by about thirty points and answer
different questions.

Uncertainty is quoted across \emph{seeds} for claims about a method, since that
is what varies when the method is rerun. For comparisons against a fixed
baseline measured on the same problems, I use a paired bootstrap over
environments instead, for reasons Sec.~\ref{sec:results} makes concrete.

\begin{table}[t]
\centering\small
\begin{tabular}{lrrr}
\toprule
Planner & Free (\%) & Clearance & s/query \\
\midrule
straight line                   & 15.6 & 0.055 & $<$0.001 \\
TrajOpt                         & 74.0 & 0.026 & 0.352 \\
CHOMP                           & 88.8 & 0.019 & 0.394 \\
RRT-Connect                     & 99.6 & 0.017 & 0.284 \\
expert (RRT+CHOMP)              & 100.0 & 0.037 & 0.342 \\
\midrule
flow, world frame               & 14.6 & $-0.071$ & 0.053 \\
flow, $(\svec,\gvec)$ reduction & \textbf{51.1} & $-0.002$ & 0.063 \\
\bottomrule
\end{tabular}
\caption{The same $500$ held-out problems solved five other ways. The learned
rows are the converged $60$-environment arms, three-seed means. Classical timings
are one CPU core and one query; the learned rows are one GPU and twenty samples,
so the time column is not like-for-like.}
\label{tab:baselines}
\end{table}

\paragraph{The floor is the finding.} A straight line from start to goal is
collision-free on $15.6\%$ of these problems. The world-frame model converges to
$14.60 \pm 0.20\%$. \emph{Per sample, whatever it has learned is worth nothing
over connecting the endpoints with a line.} I originally wrote that it was
significantly \emph{below} the line, quoting the across-seed spread. That was the
wrong denominator: the spread describes variation between training runs, not the
uncertainty of a comparison against a fixed baseline on the same problems. The
right test is paired. Resampling environments and differencing within each
replicate gives $-1.02$ points, a $95\%$ interval of $[-3.06, +0.90]$ and a
standard error of $1.01$, less than half what either rate carries alone. The
claim I can defend is that world coordinates do not \emph{beat} a straight line,
not that they lose to one.

\paragraph{Which number is the headline.} Three levels appear in this report.
$51.10 \pm 0.25$ against $14.60 \pm 0.20$ is the headline: sixty environments,
trained to a plateau, three seeds, and the setting every comparison below uses.
The $45.5$ that appears in the scaling grid is a fixed $20$-epoch budget, which
exists to compare data scales at equal optimisation. The $55.8$ at $250$
environments is the largest number I measured and rests on a single seed, which
is why it is not the headline.

\section{Results}
\label{sec:results}

\subsection{The reduction scales; the baseline does not}

\begin{table}[t]
\centering\small
\begin{tabular}{lrrrr}
\toprule
Training environments & 20 & 60 & 150 & 250 \\
\midrule
world frame           & 12.86 & 13.57 & 14.75 & 15.37 \\
$(\svec,\gvec)$ reduction & 19.23 & 34.43 & 40.60 & \textbf{45.45} \\
\midrule
gap                   & $+6.4$ & $+20.9$ & $+25.9$ & $\mathbf{+30.1}$ \\
\bottomrule
\end{tabular}
\caption{A fixed $20$-epoch budget at four data scales, three seeds per cell.
The control gains $2.6$ points across a $12.5\times$ increase in layouts; the
reduction gains $26.2$ and has not plateaued.}
\label{tab:scaling}
\end{table}

Table~\ref{tab:scaling} answers the objection I expected most: that
canonicalisation is a small-data crutch. The gap \emph{widens} monotonically
with data. Trained to convergence at sixty environments the gap is
$+36.5 \pm 0.32$ points, and the control gains nothing from a $15\times$ budget
increase --- it moves $14.6 \to 14.4$ and stays at the floor. Minimum clearance
agrees independently: the control sits at $-0.07$ throughout while the reduction
approaches feasibility at $-0.002$.

\subsection{All three mechanisms for the sixth degree of freedom are worth almost nothing}

\begin{table}[t]
\centering\small
\begin{tabular}{lrl}
\toprule
Mechanism & Worth & \\
\midrule
$(\svec,\gvec)$ reduction (5 DOF) & $+36.5$ & $\pm 0.32$, $n=3$ \\
\quad of which: 3 translations    & $+12.3$ & recentring \\
\quad of which: 2 rotations       & $+18.0$ & axis alignment \\
\midrule
roll augmentation (data)          & $+0.8$  & $\pm 0.3$ paired, n.s. \\
frame averaging (operator)        & $+0.68$ & $\pm 0.13$, $n=3$ \\
$\SO(2)$-equivariant net (weights)& $+0.15$ & seed-matched, $n=3$ \\
\bottomrule
\end{tabular}
\caption{The five removable degrees of freedom against the three mechanisms
built for the sixth.}
\label{tab:mechanisms}
\end{table}

Three unrelated ways of imposing the same symmetry all recover under a point.
That is a considerably stronger statement than any one of them makes alone, and
it is the result I find most satisfying, because for most of the internship I
had only the first two and was arguing from their flatness.

\subsection{The third mechanism: I built the equivariant architecture}
\label{sec:arch}

I had argued from the smallness of $r$ that a constrained architecture would gain
little. Sec.~\ref{sec:twosigns} shows that inference is unsound, so I built it.

The backbone is exactly $\SO(2)$-equivariant about the start--goal axis by
construction. Features are separated by irreducible representation --- the
component along the axis is invariant, the pair in the normal plane rotates
together as one complex number --- and only operations respecting that separation
are allowed: complex-linear channel mixing, obstacle pooling with attention
weights computed from invariants alone, normalisation of the rotating stream by
an invariant, and gating of that stream by scalars. Equivariance is verified on
an \emph{untrained} network at $1.2 \times 10^{-6}$, and survives multiplying
every parameter by three. I did not match capacity downward: the constrained
model carries $2{,}582{,}233$ parameters against $2{,}161{,}283$, or $19.5\%$
more, because a constrained model that is also smaller confounds architecture
with capacity.

\begin{table}[t]
\centering\small
\begin{tabular}{lrrr}
\toprule
Epoch & equivariant & unconstrained & $\Delta$ \\
\midrule
$10$  & 35.39 & 22.25 & $+13.14$ \\
$40$  & 44.91 & 38.17 & $+6.74$ \\
$100$ & 48.82 & 45.98 & $+2.84$ \\
$300$ & 51.27 & 51.12 & $+0.15$ \\
\bottomrule
\end{tabular}
\caption{Sixty environments, identical data and schedule, seed-matched. The
constraint buys training efficiency and not a higher ceiling.}
\label{tab:arch}
\end{table}

The answer is that the constraint buys optimisation rather than accuracy. It
reaches by epoch $10$ a level the unconstrained arm needs about $32$ epochs for,
and by epoch $40$ one it needs about $87$: a saving of two to three times,
largest early and shrinking as both converge. At convergence the difference is
$+0.15$ points, which is zero.

\paragraph{A bug worth recording.} My first implementation plateaued at $22.2\%$
and the mathematics was correct throughout --- every equivariance test passed at
$10^{-6}$. The fault was in how information moved through it. I was pooling the
obstacles' rotating features by taking their mean, and the mean of a roughly
symmetric cloud of directions is close to nothing: I measured that it retained
$14\%$ of a typical obstacle's magnitude. The one channel carrying ``which way do
the obstacles lie'' was arriving empty. \emph{Correct and useless are different
failure modes, and only one of them has a test.}

\subsection{What the residual does not predict}
\label{sec:twosigns}

\begin{table}[t]
\centering\small
\begin{tabular}{lrr}
\toprule
Arm ($250$ environments) & $r$ under $\SE$ & Free (\%) \\
\midrule
world frame                  & 0.0881 & 15.4 \\
\quad $+$ $\SE$ augmentation & 0.0181 & 17.5 \\
$(\svec,\gvec)$ reduction    & 0.0174 & \textbf{45.6} \\
\bottomrule
\end{tabular}
\caption{Augmentation and canonicalisation leave the field equally equivariant
and differ by $28$ points.}
\label{tab:twosigns}
\end{table}

Augmentation drives the residual down by a factor of $4.9$ --- it genuinely
teaches the network the symmetry, to the same degree the reduction enforces it by
construction --- and buys $2.1$ points. The reduction, at an indistinguishable
residual, buys $30.1$. So $r$ does not forecast what a mechanism is worth, and
the tempting rule ``build the constrained architecture only if $r$ is large'' has
no support. One number here deserves saying out loud: the world-frame arm
performs at the straight-line floor and is still $91\%$ equivariant. A model that
has learned nothing useful has landed close to the equivariant subspace.

\subsection{Is the benefit simply letting the encoder see the query?}

This is the sharpest objection to the whole report, and it is testable. Linear
probes show the world-frame scene code has within-environment standard deviation
$0.000000$ --- exactly zero, because the encoder sees only obstacles and cannot
vary with the query by construction --- while the reduced code varies almost
entirely within an environment. So perhaps the reduction buys query-dependence
rather than canonicalisation, and a one-line change to the encoder would buy the
same thing.

I built the arm that separates them: the world frame throughout, no change of
frame anywhere, with the raw query concatenated to every obstacle before the
pool. Its scene code does vary with the query, and it carries $0.07\%$ more
parameters. It converges to $15.59 \pm 0.38\%$ against the world frame's
$14.60 \pm 0.20$, both over three seeds; seed-matched the difference is $+1.01$
and positive on every seed. \emph{Query-dependence alone is worth about one point
of the $36.5$; canonicalisation supplies the other $35.5$.}

Where it lands is the part I find most useful. At $15.59\%$ it sits on the
straight line's $15.6\%$ rather than above it. Giving the encoder the query is
exactly enough to reach the trivial baseline and not enough to beat it. What
takes the model past the floor is the change of frame.

\subsection{Local geometry, and how much of my result it explains}
\label{sec:localgeom}

Late in the internship I ran an experiment I expected to be a diagnostic and
which turned out to be the largest effect in the project. The obstacle encoder
pools forty obstacles into one vector applied identically to all sixty-four
waypoints, so an individual waypoint has no channel through which to ask what is
near \emph{it} --- although collision avoidance is precisely that question. I
appended to each waypoint the signed distance to the scene and its gradient. This
is a featurisation, not an oracle: the signed distance is a deterministic
function of the obstacle primitives the network already receives. What it assumes
is that obstacles are available as explicit primitives at inference.

\begin{table}[t]
\centering\small
\begin{tabular}{lrrr}
\toprule
$60$ envs, converged & global only & $+$ local & geometry adds \\
\midrule
world frame                    & 14.60 & 54.54 & $+40.0$ \\
$(\svec,\gvec)$ reduction      & 51.10 & 79.37 & $+28.3$ \\
\quad $+$ equivariant backbone & 51.27 & \textbf{82.13} & $+30.9$ \\
\midrule
\emph{reduction contributes}   & $+36.5$ & $+24.8$ & \\
\bottomrule
\end{tabular}
\caption{Every cell is a mean over two or three seeds. The local-geometry arms
are the least seed-sensitive measurements in this work.}
\label{tab:localgeom}
\end{table}

The uncomfortable number first. \textbf{Local geometry is worth more than my
reduction is}, $+40.0$ against $+36.5$. The representation change this report is
about is not the largest effect available on its own benchmark, and I would
rather say that myself than have it pointed out to me.

The obvious follow-up is whether the reduction was only ever a workaround for a
weak encoder. If it were, supplying per-waypoint geometry should absorb its
benefit. \textbf{It does not}: with local geometry present the reduction is still
worth $+24.8$ points. I ran that cell specifically to find out whether it would
overturn my main result, and it did not. It does decay with data --- comparing at
a matched budget and differencing within seed, the reduction alongside local
geometry is worth $+22.37 \pm 0.46$ at sixty environments and
$+17.24 \pm 0.17$ at $250$ --- so anyone extrapolating to a much larger dataset
should expect this margin to narrow rather than hold.

One last observation, which I like because the architecture predicts it: the
equivariant backbone is worth $+0.15$ points without local geometry and $+2.8$
with it. The rotating stream has nothing to act on until directional information
reaches individual waypoints.

\subsection{Three shorter results}

\textbf{Two bottlenecks, not one.} Tripling the training budget is worth $+8.0$
points to the canonicalised arm and $+0.9$ to the world-frame one, an interaction
of $+7.1$. Compute cannot buy what the representation forecloses. On the
best-of-$20$ metric the same tripling moves the world-frame arm \emph{backwards},
from $42.3\%$ to $39.2\%$: more training bought a per-sample gain while costing
sample diversity.

\textbf{It replicates in a second state space.} I built an SE(3) rigid-body
domain, an L-shaped union of five spheres whose state is a full pose rather than
a point, deliberately chosen to have trivial symmetry so the domain is not
secretly easier. The world-frame arm fails to clear its own trivial floor in both
domains, the reduction clears it in both, and the ratio between arms is
$2.9\times$ in each.

\textbf{It is not specific to flow matching.} Replacing the objective with DDPM
and holding everything else fixed reproduces the gap: $+37.0$ at convergence
against flow matching's $+36.5$. The two objectives disagree early and agree at
convergence.

\section{Corrections}
\label{sec:corrections}

Five claims I made during the internship turned out to be wrong. I record them
because in each case the corrected statement is more useful than the original.

\paragraph{1. I thought the reduced representation was bit-identical under rigid
motion.} It is not; it is invariant up to one roll. The correction is what
motivates half of this report.

\paragraph{2. I reported roll augmentation as worth $+4.4$ points.} An artefact
of comparing a three-seed mean against a single-seed ablation. Seed-matched it is
$+0.8 \pm 0.3$ and not significant. The general lesson: the across-seed spread is
five times the paired spread, so any ablation compared unpaired at this scale is
measuring seed noise.

\paragraph{3. I argued the roll gauge could be fixed by a $C_4$ argument.} That
holds only when the axis is a four-fold axis of the cube. The correct hypothesis
is central inversion, which is weaker, true, and holds for a much broader class
of obstacle distributions than my own generator.

\paragraph{4. I reported the residual as falling with scale.} That came from two
points under an older definition. Measured at four scales under the definition
that makes the bound an identity, it is flat. I had built an argument on the
trend, and that argument does not survive.

\paragraph{5. I decided not to build the equivariant architecture, and gave two
reasons, both wrong.} I argued that $r$ was small and falling, so the
architecture would enforce something arriving on its own. The trend was an
artefact and the level does not predict. When I finally built it the ceiling was
indeed unchanged --- so I had reached the right answer by an argument that does
not support it, and missed entirely the two-to-threefold saving in optimisation.
This is the correction I find most instructive.

\section{Discussion and limitations}
\label{sec:limits}

\paragraph{Which robots this transfers to.} The reduction needs the group to act
on the \emph{state space}, not just the workspace. That holds for a point mass
and for the rigid body of the second domain. It fails for a manipulator planning
in joint space: the query still defines a frame in the workspace, but a rigid
motion of the world induces no clean action on joint angles. Mobile bases,
free-flying bodies and end-effector-space planning pass this test; joint-space
planning for arms does not.

\paragraph{The learned planner loses to the classical ones.} RRT-Connect solves
$99.6\%$ in under $0.3$\,s on one CPU core; the learned sampler reaches $51.1\%$.
On this benchmark one should use RRT-Connect. The contribution is a controlled
measurement of what a representation change is worth, not a planner proposal, and
that measurement is only isolable where the confounds are removable.

\paragraph{Scope.} One workspace family, one clutter generator, one architecture.
The claim that the sixth degree of freedom needs no treatment relies on the
training distribution supplying implicit roll diversity through many start--goal
directions; a dataset with few directions per environment could plausibly reverse
it.

\paragraph{Seeds.} The $60$-environment cells carry three seeds, the
local-geometry arms two or three, and the $250$-environment cells one. Every
claim I lean on has at least two.

\section{What I would do next}

\textbf{Make $r$ vary deliberately.} Training on a narrow cone of start--goal
directions should starve the model of implicit roll diversity and drive $r$ up.
If frame averaging then starts to pay in proportion, $r$ becomes a decision rule
for post-hoc symmetrisation; if it does not pay even at large $r$, that is the
harder negative. I attempted this and the attempt failed: narrowing the cone
raises $r$ but degrades quality at the same time, and I could not separate them.

\textbf{Capacity as a controlled variable.} The model is small, and "the model is
small" is currently a limitation rather than something I measured. Doubling the
trunk on both arms would turn it into a controlled variable.

\textbf{More environments.} The scaling curve has not plateaued at $250$, which
is the dataset cap. Cutting paths per pair from thirty to five would buy roughly
six times the environments for the same generation cost.

\section{Conclusion}

Of the six degrees of freedom of $\SE$, five are removable exactly and for free
by a frame the query already supplies, and that is worth $+36.5$ points on this
benchmark against a world-frame arm that does not beat a straight line. The sixth
is not removable by any continuous rule, and all three mechanisms available for
it --- augmentation, frame averaging, and an equivariant architecture --- are
worth under a point each. The architecture is not pointless: it reaches any given
level two to three times sooner, so the symmetry buys optimisation rather than
accuracy.

I also bound my own contribution from above. Supplying per-waypoint geometry is
worth more than the frame is, and the frame still adds $+24.8$ points alongside
it. The recommendation I would make for this benchmark is correspondingly
narrow: canonicalise the degrees of freedom that admit a continuous canonical
form, because the query usually supplies more of them than one expects, and
measure the rest rather than assuming them.

\end{document}
